\subsubsection{Reward A: Format Compliance}

\paragraph{Definition.}
Format compliance is a binary reward that returns 1.0 if the model's response contains a parseable answer in the expected \texttt{\#\#\#\# <number>} format, and 0.0 otherwise.
Formally:
\[
R_A(\text{response}) =
\begin{cases}
1.0 & \text{if } \texttt{\#\#\#\# (-?[0-9.,]+)} \text{ matches} \\
0.0 & \text{otherwise}
\end{cases}
\]

\paragraph{Rationale.}
Analysis of the SFT baseline (15.39\% strict accuracy) reveals that approximately 11\% of responses fail to produce any parseable answer.
These responses receive zero reward from the correctness function regardless of their reasoning quality, providing no learning signal.
By rewarding format compliance independently, we decouple format learning from mathematical reasoning.
This is particularly valuable early in training when most or all 16 sampled responses per problem may be incorrect: the format reward still provides non-zero advantage variance, enabling gradient updates even when correctness-based reward is uniformly zero.

\paragraph{Literature.}
Format compliance rewards are standard practice in RLVR implementations for GSM8K.
The GSM8K-RLVR benchmark, veRL framework, and Unsloth RL documentation all include format-checking as a reward component.
The motivation is consistent across these works: ensuring the model learns to produce extractable answers is a prerequisite for the correctness signal to be effective.

\paragraph{Alternatives considered.}
We considered a calculator usage reward (rewarding invocation of calculator tools) and a reasoning structure reward (rewarding chain-of-thought formatting).
Calculator usage was rejected because it risks incentivizing meaningless tool calls that do not contribute to problem solving.
Reasoning structure was rejected because it is susceptible to reward hacking via padding attacks, where the model inflates response length without improving reasoning quality (Gao et al., 2024).
